\section{问题一：典型风光场景下运行指标分析}

\subsection{问题分析与求解思路}

问题一要求在电解槽与合成氨装置每日满负荷连续运行条件下，计算园区典型日的功率平衡关系和各项运行指标。由于所有设备均满负荷运行，不存在决策变量，本质上是一个确定性的功率平衡计算问题。求解思路为：逐时段计算净负荷，确定购售电功率，累加得到日能量指标，最终评估绿电直连三项指标的满足情况。

\begin{figure}[H]
\centering
\includegraphics[width=0.85\textwidth,height=0.38\textheight,keepaspectratio]{../figures/fig_flow_q1.pdf}
\caption{问题一求解流程图}\label{fig:flow_q1}
\end{figure}

\subsection{数学模型建立}

\subsubsection{功率平衡方程}

在每个时段$t \in \{1, 2, \ldots, 24\}$，园区功率平衡关系为：
\begin{equation}\label{eq:power_balance_q1}
P_{wind}(t) + P_{pv}(t) + P_{buy}(t) = P_{load}(t) + P_{H2NH3}^{max} + P_{sell}(t)
\end{equation}

由于满负荷运行，制氢氨设备总功率为常数：
\begin{equation}
P_{H2NH3}^{max} = P_{ALKEL}^{max} + P_{PEMEL}^{max} + P_{NH3}^{max} = 10 + 10 + 0.75 = 20.75 \text{ MW}
\end{equation}

定义净负荷为园区总用电需求减去新能源发电：
\begin{equation}
P_{net}(t) = P_{load}(t) + P_{H2NH3}^{max} - P_{wind}(t) - P_{pv}(t)
\end{equation}

根据购售电互斥原则，同一时段不能同时购电和售电：
\begin{equation}
P_{buy}(t) = \max(0, P_{net}(t)), \quad P_{sell}(t) = \max(0, -P_{net}(t))
\end{equation}

\subsubsection{日能量指标计算}

日总用电量、新能源发电量、网购电量和上网电量分别为：
\begin{equation}
E_{total} = \sum_{t=1}^{24} [P_{load}(t) + P_{H2NH3}^{max}] \cdot \Delta t
\end{equation}
\begin{equation}
E_{RE} = \sum_{t=1}^{24} [P_{wind}(t) + P_{pv}(t)] \cdot \Delta t
\end{equation}
\begin{equation}
E_{buy} = \sum_{t=1}^{24} P_{buy}(t) \cdot \Delta t, \quad E_{sell} = \sum_{t=1}^{24} P_{sell}(t) \cdot \Delta t
\end{equation}

\subsubsection{绿电直连指标}

根据政策文件定义，三项绿电直连指标的计算公式为：
\begin{equation}\label{eq:R_self}
R_{self} = \frac{E_{total} - E_{sell} - E_{buy}}{E_{RE}} \quad (\text{要求} > 60\%)
\end{equation}
\begin{equation}\label{eq:R_green}
R_{green} = \frac{E_{RE} - E_{sell}}{E_{total}} \quad (\text{要求} > 30\%)
\end{equation}
\begin{equation}\label{eq:R_grid}
R_{grid} = \frac{E_{sell}}{E_{RE}} \quad (\text{要求} < 20\%)
\end{equation}

\subsubsection{吨氨成本模型}

吨氨成本由新能源度电成本、购电成本、设备运维成本和售电收入四部分构成：
\begin{equation}\label{eq:C_ton}
C_{ton} = \frac{C_{RE} + C_{buy} + C_{OM} - I_{sell}}{Q_{NH3}}
\end{equation}

其中各分项成本为：
\begin{equation}
C_{RE} = \sum_{t=1}^{24} [c_{wind} \cdot P_{wind}(t) + c_{pv} \cdot P_{pv}(t)] \cdot \Delta t
\end{equation}
\begin{equation}
C_{buy} = \sum_{t=1}^{24} \lambda_{buy}(t) \cdot P_{buy}(t) \cdot \Delta t
\end{equation}
\begin{equation}
C_{OM} = \sum_{t=1}^{24} [c_{ALKEL} \cdot P_{ALKEL}^{max} + c_{PEMEL} \cdot P_{PEMEL}^{max} + c_{NH3} \cdot P_{NH3}^{max}] \cdot \Delta t
\end{equation}
\begin{equation}
I_{sell} = \lambda_{sell} \cdot E_{sell}
\end{equation}

分时电价按峰平谷三段划分：低谷时段（23:00--07:00）电价0.3424元/kWh，平时段（07:00--10:00、15:00--18:00、21:00--23:00）电价0.6074元/kWh，高峰时段（10:00--15:00、18:00--21:00）电价0.8024元/kWh。

\subsection{计算结果与分析}

\subsubsection{功率平衡曲线}

基于附件1的常规负荷标幺功率曲线和附件2的典型日风光标幺功率曲线，计算得到园区典型日24小时功率平衡结果（图\ref{fig:q1_power_balance}）。

\begin{figure}[H]
\centering
\includegraphics[width=0.9\textwidth,height=0.38\textheight,keepaspectratio]{../figures/fig_q1_power_balance.pdf}
\caption{典型日功率平衡曲线}\label{fig:q1_power_balance}
\end{figure}

典型日功率平衡曲线呈现明显的时序特征：夜间（0:00--7:00）风电出力较高但光伏为零，园区净负荷为正，需要从电网购电；白天（9:00--16:00）光伏出力达到峰值，叠加风电后新能源总出力远超园区用电需求，产生大量余电上网；傍晚至夜间（17:00--24:00）光伏出力迅速下降，园区再次转为购电状态。这种"白天余电、夜间缺电"的时序不匹配特征是绿电直连指标不达标的根本原因。

\subsubsection{运行指标计算结果}

根据功率平衡计算，园区典型日各项能量指标和绿电直连指标如表\ref{tab:q1_results}所示。

\begin{table}[H]
\centering
\caption{问题一典型日运行指标计算结果}\label{tab:q1_results}
\begin{tabular}{lccc}
\toprule
指标 & 计算值 & 政策要求 & 是否满足 \\
\midrule
日总用电量 $E_{total}$ & 558.72 MWh & --- & --- \\
新能源发电量 $E_{RE}$ & 603.45 MWh & --- & --- \\
网购电量 $E_{buy}$ & 172.04 MWh & --- & --- \\
上网电量 $E_{sell}$ & 216.77 MWh & --- & --- \\
自发自用比例 $R_{self}$ & 28.16\% & $>60\%$ & 不满足 \\
绿电比例 $R_{green}$ & 69.21\% & $>30\%$ & 满足 \\
上网电量比例 $R_{grid}$ & 35.92\% & $<20\%$ & 不满足 \\
吨氨成本 $C_{ton}$ & 4322.34 元/吨 & --- & --- \\
\bottomrule
\end{tabular}
\end{table}

计算结果表明，尽管新能源总发电量（603.45 MWh）大于日总用电量（558.72 MWh），园区在能量总量上具备自给能力，但三项绿电直连指标中仅绿电比例（69.21\%）满足要求。自发自用比例仅为28.16\%，远低于60\%的政策要求；上网电量比例高达35.92\%，远超20\%的上限。

\begin{figure}[H]
\centering
\includegraphics[width=0.85\textwidth,height=0.38\textheight,keepaspectratio]{../figures/fig_q1_indicators.pdf}
\caption{绿电直连三项指标实际值与阈值对比}\label{fig:q1_indicators}
\end{figure}

指标不达标的根本原因在于风光出力与用电负荷的时序不匹配（图\ref{fig:q1_indicators}）。园区制氢氨设备24小时满负荷运行，用电需求恒定，而光伏出力集中在白天6--8小时，导致白天大量余电上网（216.77 MWh），夜间则需大量购电（172.04 MWh）。这种"发用时序错配"使得自发自用比例偏低、上网比例偏高。要改善绿电指标，需要通过调整制氢氨设备的运行时段（使其与风光出力时序匹配）或配置储能来实现"削峰填谷"。

\subsubsection{吨氨成本分解}

典型日吨氨成本为4322.34元/吨，其成本构成为：新能源度电成本79765.20元（占51.3\%）、购电成本97721.03元（占62.8\%）、设备运维成本60036.00元（占38.6\%），扣除售电收入81918.06元（占$-$52.6\%）。购电成本是最大的成本项，主要来源于夜间高峰时段的购电需求，这进一步说明了优化制氢氨设备运行时段的必要性。
